\documentclass[11pt,a4paper]{article}

\usepackage{amsmath,amssymb,bm}
\usepackage{geometry}
\usepackage{hyperref}
\usepackage{graphicx}
\usepackage{booktabs}
\usepackage{xcolor}

\geometry{margin=2.5cm}

\newcommand{\be}{\hat{\mathbf{e}}}
\newcommand{\bn}{\hat{\mathbf{n}}}
\newcommand{\bl}{\hat{\mathbf{l}}}
\newcommand{\bc}{\hat{\mathbf{c}}}
\newcommand{\bff}{\mathbf{f}}
\newcommand{\bx}{\mathbf{x}}
\newcommand{\R}{\mathbb{R}}

\title{Left Ventricle Surface Geometry and Helical Fiber Architecture\\
\large Methodology and Formulation}
\author{BioFEM-IA}
\date{\today}

\begin{document}
\maketitle
\tableofcontents
\vspace{1em}

% -----------------------------------------------------------------------
\section{Overview}

This document describes the analytical geometry and rule-based myofiber
architecture used to construct a simplified left-ventricle (LV) surface mesh
in BioFEM-IA (\texttt{src/lv\_geometry.c}).  The model is intended as a
lightweight starting point: no volumetric or imaging data are required, and
the mesh is generated entirely from closed-form expressions.

The mesh is a \emph{surface} (shell) representation of the
\textbf{midmyocardial layer} of the LV free wall, suitable for thin-shell
finite-element analyses.  Fiber directions are assigned element-by-element via
a classical rule-based interpolation of Streeter's experimental midwall
measurements~\cite{Streeter1969}, following the algorithmic framework of Bayer
et al.~\cite{Bayer2012}.

% -----------------------------------------------------------------------
\section{Left Ventricle Geometry}

\subsection{Prolate Spheroid Parameterisation}

The left ventricle is approximated as a \emph{truncated prolate spheroid}
(oblate-spheroid prolate ellipsoid of revolution).  In the reference
configuration, the long axis coincides with the $z$-axis.  The apex is
placed at the top of the long axis, $(0,0,a)$, and the open base rim lies in
a plane of constant $z$ below the equator.

The surface is parameterised by two angles,
\begin{equation}
  \bx(\theta,\varphi)
  =
  \begin{pmatrix}
    b\sin\theta\cos\varphi \\
    b\sin\theta\sin\varphi \\
    a\cos\theta
  \end{pmatrix},
  \qquad
  \theta\in[0,\theta_{\max}],\quad
  \varphi\in[0,2\pi),
  \label{eq:parameterisation}
\end{equation}
where
\begin{itemize}
  \item $a > 0$ is the \textbf{polar} (long-axis) semi-axis,
  \item $b > 0$ is the \textbf{equatorial} semi-axis ($b < a$ for a prolate
        shape),
  \item $\theta$ is the \textbf{polar angle} measured from the apex
        ($\theta=0$ at the apex, increasing toward the base), and
  \item $\varphi$ is the \textbf{azimuthal angle}.
\end{itemize}

The level-set function of the spheroid is
\begin{equation}
  F(\bx) = \frac{x^2}{b^2} + \frac{y^2}{b^2} + \frac{z^2}{a^2} - 1 = 0.
\end{equation}

\subsection{Truncation Height}

A single dimensionless parameter $h \in (0,1]$, called the \emph{cut
fraction}, controls how much of the spheroid is retained.  It is defined as
the fraction of the full axial height $2a$ kept from the apex:
\begin{equation}
  h = \frac{a - z_{\text{base}}}{2a}
  \;\Longrightarrow\;
  z_{\text{base}} = a(1 - 2h).
\end{equation}
The corresponding maximum polar angle is
\begin{equation}
  \boxed{\theta_{\max} = \arccos(1 - 2h).}
  \label{eq:theta_cut}
\end{equation}
\begin{itemize}
  \item $h = 0.5$: hemisphere ($z_{\text{base}} = 0$, equatorial cut).
  \item $h = 0.65$: typical LV truncation ($z_{\text{base}} < 0$, base extends
        below the equator).
  \item $h = 1.0$: full spheroid (no truncation).
\end{itemize}

\subsection{Discrete Mesh}

The surface is discretised with a structured triangulation.

\subsubsection*{Node layout}

\begin{itemize}
  \item \textbf{Node 0} (apex): $\bx = (0, 0, a)$.
  \item \textbf{Ring $k$} ($k = 0, 1, \dots, N_\theta - 1$):
        $N_\varphi$ nodes uniformly spaced in $\varphi$, at polar angle
        \[
          \theta_k = (k+1)\,\frac{\theta_{\max}}{N_\theta}.
        \]
        Node index: $1 + k\,N_\varphi + j$ for azimuthal position
        $j = 0, \dots, N_\varphi - 1$.
\end{itemize}

Total nodes: $N_v = 1 + N_\theta N_\varphi$.

\subsubsection*{Element connectivity}

\begin{enumerate}
  \item \textbf{Apex fan} ($N_\varphi$ triangles): connect the apex to
        consecutive pairs of ring-0 nodes,
        \[
          T_j^{\rm apex} = \{0,\;1+j,\;1+(j+1)\bmod N_\varphi\},
          \quad j = 0,\dots,N_\varphi-1.
        \]
  \item \textbf{Latitude bands} ($N_\theta - 1$ bands, each with
        $2N_\varphi$ triangles): each latitude quad is split along its
        diagonal into two triangles.  For band $k$ (between ring $k$ and
        ring $k+1$) and column $j$:
        \begin{align*}
          T_1 &= \{n_{bl},\;n_{tl},\;n_{br}\}, \\
          T_2 &= \{n_{tl},\;n_{tr},\;n_{br}\},
        \end{align*}
        where $n_{bl/br/tl/tr}$ are the bottom-left, bottom-right, top-left,
        and top-right corners of the quad.
\end{enumerate}

Total elements: $N_e = N_\varphi(2N_\theta - 1)$.

\subsubsection*{Patch nodes ($nv4/nv5/nv6$)}

For subdivision-surface bending (curvature mode~1), each element requires the
``opposite'' node of each of its three edge-sharing neighbours.  These are
computed with an $O(N_e^2)$ neighbour search (acceptable for the small meshes
used in this module) and stored in \texttt{ibm->nv4/nv5/nv6}.  Elements on
the open base boundary use the sentinel value $10^6$ to indicate no neighbour.

% -----------------------------------------------------------------------
\section{Helical Myofiber Architecture}

\subsection{Myocardial Wall Layers}

The LV free wall consists of three concentric layers.

\begin{itemize}
  \item \textbf{Endocardium} — thin inner endothelial lining of the ventricular
        chamber.  Passively transmits luminal pressure; not contractile.
  \item \textbf{Myocardium} — the thick muscular middle layer, composed almost
        entirely of cardiomyocytes whose sarcomeres generate active tension.
        This is the layer responsible for the pumping function of the heart.
        Within the myocardium, the \emph{subendocardial} portion (inner third)
        is activated first via the Purkinje network, but contractile force is
        generated throughout the full thickness.
  \item \textbf{Epicardium} — thin outer serous membrane (visceral pericardium).
        Passively supports the wall; not contractile.
\end{itemize}

The myocardium is therefore the layer of primary interest for muscle-activation
models.

\subsection{Transmural Fiber Architecture and the Representative Surface}

Streeter et al.~\cite{Streeter1969} measured the myofiber helix angle
$\alpha$ across the wall thickness (transmurally) at multiple apex-to-base
stations.  The angle varies \emph{continuously} from a large positive value at
the endocardium to a large negative value at the epicardium, passing through
approximately zero at the midwall:

\begin{table}[h]
\centering
\begin{tabular}{lll}
\toprule
Layer & Depth & Typical helix angle $\alpha$ \\
\midrule
Endocardium (inner) & 0\%   & $+60^\circ$ to $+90^\circ$ \\
\textbf{Midmyocardium} & \textbf{50\%} & \boldmath$\pm 60^\circ$ \textbf{(apex $\to$ base)} \\
Epicardium (outer)  & 100\% & $-60^\circ$ to $-90^\circ$ \\
\bottomrule
\end{tabular}
\caption{Transmural helix-angle distribution in the LV wall
         (Streeter 1969, Table~1).  Sign convention: positive angles indicate
         left-handed helices (endocardial sense), negative angles right-handed
         helices (epicardial sense).}
\label{tab:transmural}
\end{table}

Because BioFEM-IA uses a \emph{single-layer} thin-shell formulation, the mesh
represents \textbf{one} depth through the wall.  The conventional choice in the
computational cardiac mechanics literature
(Bayer et al.~\cite{Bayer2012}, Bovendeerd et al.~\cite{Bovendeerd1994})
is to place that surface at the \textbf{midmyocardium} (mid-wall, 50\% depth).
This choice is motivated by two reasons:
\begin{enumerate}
  \item The midwall is the \emph{neutral surface} for bending of the shell,
        minimising the error introduced by collapsing finite thickness onto a
        single surface.
  \item The midwall helix angles ($+60^\circ$ at the apex, crossing $0^\circ$
        near the equator, $-60^\circ$ at the base) are well characterised by
        Streeter's direct measurements and serve as the canonical reference
        for rule-based fiber assignment.
\end{enumerate}

\noindent
\textbf{Summary:} the present model represents the \emph{midmyocardial surface}
of the LV free wall.  The default angle parameters
$\alpha_{\rm apex} = +60^\circ$, $\alpha_{\rm base} = -60^\circ$ are taken
directly from Streeter's midwall data~\cite{Streeter1969}.

\subsection{Motivation}

Myocardial fibers in the LV wall follow a helical path that rotates
continuously from the endocardial surface (positive helix angle) to the
epicardial surface (negative helix angle), with a smooth zero crossing at the
midwall~\cite{Streeter1969,Holzapfel2009}.  For a \emph{surface} mesh
representing the midmyocardial layer (Section~3.2), the helical rotation
manifests as a spatial variation of the fiber angle with position along the
long axis: near the apex the fibers are steeply inclined, while near the base
they are nearly circumferential.

\subsection{Local Orthonormal Frame}

At the centroid $\bx_c$ of each triangular element, three orthonormal vectors
are constructed.

\paragraph{Outward surface normal $\bn$.}
The outward unit normal to the spheroid is proportional to $\nabla F$:
\begin{equation}
  \bn = \mathrm{unit}\!\left(\frac{x_c}{b^2},\;\frac{y_c}{b^2},\;
                               \frac{z_c}{a^2}\right).
  \label{eq:normal}
\end{equation}

\paragraph{Meridional (apex-to-base) direction $\bl$.}
The apex-to-base direction is obtained by projecting the downward unit vector
$\hat{\mathbf{z}}_{-} = (0,0,-1)$ onto the tangent plane:
\begin{equation}
  \bl_{\rm raw} = \hat{\mathbf{z}}_{-}
                  - \bigl(\hat{\mathbf{z}}_{-}\cdot\bn\bigr)\,\bn,
  \qquad
  \bl = \mathrm{unit}(\bl_{\rm raw}).
  \label{eq:meridional}
\end{equation}
At the apex ($\theta\to 0$) the projection degenerates; in that case a
fallback tangent is derived by projecting $(1,0,0)$ instead.

\paragraph{Circumferential direction $\bc$.}
The right-hand circumferential direction is
\begin{equation}
  \bc = \bn \times \bl.
  \label{eq:circumferential}
\end{equation}
By construction $\{\bn,\bl,\bc\}$ is an orthonormal triad, and $\bl$, $\bc$
lie in the tangent plane of the spheroid at $\bx_c$.

\subsection{Helix Angle Model}

Following Streeter~\cite{Streeter1969} and the rule-based framework of
Bayer et al.~\cite{Bayer2012}, the helix angle $\alpha$ (measured from the
circumferential direction $\bc$ toward the meridional direction $\bl$) is
interpolated linearly along the long axis:
\begin{equation}
  \boxed{
    \alpha(z^*) = (1 - z^*)\,\alpha_{\rm apex} + z^*\,\alpha_{\rm base},
  }
  \label{eq:helix_angle}
\end{equation}
where the normalised axial coordinate is
\begin{equation}
  z^* = \frac{a - z_c}{a - z_{\rm base}} \;\in\; [0,1],
  \label{eq:zstar}
\end{equation}
with $z^* = 0$ at the apex and $z^* = 1$ at the base.

Default values (following Streeter's midwall measurements~\cite{Streeter1969}):
\[
  \alpha_{\rm apex} = +60^\circ, \qquad \alpha_{\rm base} = -60^\circ.
\]

\subsection{Fiber Unit Vector}

The fiber direction at element $e$ is the unit vector in the tangent plane
making angle $\alpha$ with $\bc$:
\begin{equation}
  \boxed{
    \bff_e = \cos\alpha\,\bc + \sin\alpha\,\bl,
  }
  \label{eq:fiber_vec}
\end{equation}
stored as \texttt{ibm->n\_fib[e]}.  Since $\bc$ and $\bl$ are both unit
vectors and mutually orthogonal, $\bff_e$ is automatically a unit vector.

\paragraph{Interpretation of $\alpha$.}
\begin{itemize}
  \item $\alpha = 0^\circ$: fiber is circumferential ($\bff = \bc$).
  \item $\alpha = 90^\circ$: fiber is meridional/longitudinal ($\bff = \bl$).
  \item $\alpha = 60^\circ$ (apex): fiber is predominantly meridional with a
        helical component.
  \item $\alpha = -60^\circ$ (base): fiber is predominantly circumferential
        but wound in the opposite helical sense.
\end{itemize}

\subsection{Variation Along the Wall}

Figure~\ref{fig:fiber_schematic} illustrates the expected fiber rotation.
The helix angle transitions smoothly from $+\alpha_{\rm apex}$ at the apex to
$-\alpha_{\rm base}$ at the base, crossing $0^\circ$ at the midpoint
$z^* = \alpha_{\rm apex}/(\alpha_{\rm apex} - \alpha_{\rm base})$ (at
$z^* = 0.5$ for the symmetric default values).

\begin{figure}[h]
\centering
% Placeholder — replace with actual figure once generated from ParaView.
\fbox{\parbox{0.7\textwidth}{\centering\vspace{2cm}
  \textit{ParaView rendering of the LV mesh with fiber vectors coloured by
  helix angle $\alpha(z^*)$. Generated from \texttt{lv\_fiber\_00.vtk}.}
  \vspace{2cm}}}
\caption{Helical fiber architecture on the truncated prolate-spheroid LV
surface. Colours indicate $\alpha$ from $-60^\circ$ (base, blue) to
$+60^\circ$ (apex, red).}
\label{fig:fiber_schematic}
\end{figure}

% -----------------------------------------------------------------------
\section{Implementation Parameters}

\begin{table}[h]
\centering
\begin{tabular}{llll}
\toprule
PETSc option & Variable & Default & Description \\
\midrule
\texttt{-lv\_geom\_process} & \texttt{lv\_geom\_process} & 0        & Enable LV mesh (0 = file input) \\
\texttt{-lv\_a}             & \texttt{lv\_a}             & 4.5 cm   & Polar semi-axis \\
\texttt{-lv\_b}             & \texttt{lv\_b}             & 2.5 cm   & Equatorial semi-axis \\
\texttt{-lv\_f\_cut}        & \texttt{lv\_f\_cut}        & 0.55     & Cut fraction $h$ \\
\texttt{-lv\_N\_theta}      & \texttt{lv\_N\_theta}      & 16       & Number of latitude rings \\
\texttt{-lv\_N\_phi}        & \texttt{lv\_N\_phi}        & 32       & Nodes per ring \\
\texttt{-lv\_alpha\_apex}   & \texttt{lv\_alpha\_apex}   & $+60^\circ$ & Apex helix angle \\
\texttt{-lv\_alpha\_base}   & \texttt{lv\_alpha\_base}   & $-60^\circ$ & Base helix angle \\
\bottomrule
\end{tabular}
\caption{PETSc command-line options for the LV geometry module.}
\label{tab:options}
\end{table}

Example invocation (from the case directory, which must contain \texttt{control.dat}):
\begin{verbatim}
cd BioFEM-studies/lv_test
biofem run --backend dmplex --np 1          # local / interactive
biofem run --backend dmplex --np 8 \        # submit to SLURM
           --account <account> --time 00:30:00
\end{verbatim}
All \texttt{-lv\_*} parameters are read from \texttt{control.dat} via PETSc's options
database; no flags need to be passed on the command line unless overriding a
value from the file.

% -----------------------------------------------------------------------
\section{Output}

When \texttt{-lv\_geom\_process~1} is active, fiber vectors are written in
two places:

\paragraph{Initialisation VTK (\texttt{lv\_fiber\_00.vtk}).}
Written once at startup to \texttt{out\_dir}.  Contains the surface mesh as
\texttt{UNSTRUCTURED\_GRID} (VTK cell type~5) and a \texttt{VECTORS} cell-data
field \texttt{fiber\_direction} with $\bff_e$ for every element.

\paragraph{Time-step surface files (\texttt{surface\{ibi\}\_\{ti\}.vtk}).}
The standard \texttt{Output()} routine (controlled by \texttt{-tiout}) writes
a \texttt{VECTORS nfib float} cell-data block to each surface VTK whenever
\texttt{lv\_geom\_process}~$> 0$.  This allows the fiber field to be viewed
alongside the deformed geometry and stress/strain fields at every saved
time step.

Open either file in ParaView and apply the \emph{Glyph} filter (arrow glyphs
on \texttt{nfib} or \texttt{fiber\_direction}) to visualise the helical fiber
architecture over the LV surface.

% -----------------------------------------------------------------------
\begin{thebibliography}{9}

\bibitem{Streeter1969}
Streeter~DD~Jr, Spotnitz~HM, Patel~DJ, Ross~J~Jr, Sonnenblick~EH.
\newblock Fiber orientation in the canine left ventricle during diastole and
  systole.
\newblock \textit{Circulation Research}, 24(3):339--347, 1969.

\bibitem{Bayer2012}
Bayer~JD, Blake~RC, Plank~G, Trayanova~NA.
\newblock A novel rule-based algorithm for assigning myocardial fiber
  orientation to computational heart models.
\newblock \textit{Annals of Biomedical Engineering}, 40(10):2243--2254, 2012.

\bibitem{Holzapfel2009}
Holzapfel~GA, Ogden~RW.
\newblock Constitutive modelling of passive myocardium: a structurally based
  framework for material characterization.
\newblock \textit{Philosophical Transactions of the Royal Society A},
  367:3445--3475, 2009.

\bibitem{Bovendeerd1994}
Bovendeerd~PHM, Arts~T, Huyghe~JM, van Campen~DH, Reneman~RS.
\newblock Dependence of local left ventricular wall mechanics on myocardial
  fiber orientation: a model study.
\newblock \textit{Journal of Biomechanics}, 27(7):941--951, 1994.

\end{thebibliography}

\end{document}
